\documentclass[11pt]{article}
\usepackage[utf8]{inputenc}
\usepackage[T1]{fontenc}
\usepackage[margin=1in]{geometry}
\usepackage{amsmath,amssymb,amsfonts}
\usepackage{booktabs}
\usepackage{graphicx}
\usepackage{hyperref}
\usepackage{microtype}
\usepackage{enumitem}
\usepackage{xcolor}

\title{\textbf{METS: A Rank-Extreme Trend Sorting Algorithm with Asymptotic Significance Testing for Lightweight Multi-Series Comparison}}
\author{Max Leah\\Independent Researcher\\\texttt{1max.leah@gmail.com}}
\date{}

\begin{document}
\maketitle

\begin{abstract}
In financial analysis, meteorological monitoring, and industrial IoT, one frequently needs to rank hundreds of time series of varying length and scale by their overall trend. Existing methods fall into two categories: ordinary least squares (OLS) regression is computationally efficient but highly sensitive to outliers, while nonparametric methods such as Theil-Sen estimation are robust but computationally expensive. This paper proposes METS (Min-Max-Extreme Trend Sorter), a trend ranking algorithm based on rank-space extreme statistics. The algorithm eliminates scale differences via variance normalization, extracts the extreme position difference in rank space as a robust estimate of trend direction and strength, and combines quality assessment (monotonic contribution ratio and pseudo-$R^2$), multi-scale segmentation fusion, and an asymptotic significance test based on Spearman rank correlation. Experiments on six real-world datasets (macroeconomic indicators of two countries, corporate panel data, sea surface temperature, CO$_2$ concentration, and sunspots) show that METS achieves 75\%--100\% agreement with Mann-Kendall trend tests in significance judgment, and 75\%--100\% direction consistency, while its ranking precision is comparable to OLS in low-noise scenarios (0.913 vs.\ 0.923). The computational complexity is $O(n \log n)$. All experimental results are fully reproducible; limitations are honestly reported.
\end{abstract}

\begin{keywords}
trend ranking; rank statistics; extreme position; asymptotic test; multi-scale validation
\end{keywords}

\section{Introduction}
\subsection{Background and Motivation}
Multi-series trend comparison is a fundamental task in data-driven decision systems. For example, fund managers need to rank thousands of stocks by growth strength, meteorologists need to rank monitoring stations by warming rate, and equipment maintenance engineers need to rank sensors by degradation rate. These tasks face three common challenges: (1) a large number of series (hundreds to thousands); (2) varying lengths and scales across series; and (3) the presence of outlier noise. The core requirement is to achieve reasonable ranking with low computational cost while ensuring cross-series comparability and statistical credibility.

\subsection{Existing Methods and Limitations}
Existing methods can be categorized into three classes:

\begin{itemize}
\item \textbf{Regression-based methods}: ordinary least squares (OLS) linear regression, using the regression slope as the ranking criterion. OLS is theoretically mature and computationally efficient ($O(n)$), but is extremely sensitive to outliers---a single anomalous point can significantly alter the slope estimate. It also cannot assess whether a trend is statistically significant, and does not handle scale differences across series.
\item \textbf{Nonparametric robust estimation}: Theil-Sen estimation (Theil, 1950; Sen, 1968) takes the median of all pairwise slopes, achieving near-optimal robustness (breakdown point of 29.3\%). It is the standard tool for trend analysis in meteorology. However, its naive implementation has $O(n^2)$ complexity (all pairwise combinations), which is prohibitive for large-scale multi-series scenarios. The Mann-Kendall test (Mann, 1945; Kendall, 1948) belongs to this class but provides only significance, not a continuous ranking strength.
\item \textbf{Extreme-index methods} (Hu et al., 2003): use extreme positions to construct indices for fast searching, but target intra-series element ranking rather than inter-series trend ranking.
\end{itemize}

Zhang et al.\ (2016) proposed a slope-based non-uniform sampling method (SBNS) for function-type data, but its core still relies on curve fitting.

\subsection{Contributions and Positioning}
This paper proposes METS. We emphasize an honest, exploratory positioning: through systematic experiments, we report the algorithm's real performance and limitations rather than claiming unrealistic advantages. Our contributions are:

\begin{enumerate}
\item \textbf{Scale invariance}: variance normalization makes the slope a dimensionless shape measure;
\item \textbf{Noise-robust rank-space extrema}: extreme position difference in rank space avoids contamination by outliers in the raw value domain;
\item \textbf{Asymptotic significance testing}: a Spearman rank correlation $t$-approximation replaces the computationally expensive permutation test, reducing the full pipeline complexity to $O(n \log n)$;
\item \textbf{Systematic validation}: seven synthetic and six real-data experiments honestly assess precision, robustness, efficiency, ablation, and parameter sensitivity.
\end{enumerate}

\section{METS Algorithm Design}
\subsection{Problem Definition and Normalization}
Let there be $N$ series $\mathbf{A}_1, \mathbf{A}_2, \dots, \mathbf{A}_N$, where the $i$-th series has length $n_i$, denoted $\mathbf{A}_i = [a_{i,1}, a_{i,2}, \dots, a_{i,n_i}]$. Index $t$ represents time or sequential position.

\textbf{Scale normalization}: to make series of different scales comparable, each series is normalized by its standard deviation:
\begin{equation}
\tilde{a}_{i,t} = \frac{a_{i,t} - \mu_i}{\sigma_i}, \quad t = 1, \dots, n_i
\end{equation}
where $\mu_i$ and $\sigma_i$ are the mean and standard deviation. After normalization, each series has zero mean and unit variance, making the slope a dimensionless quantity. Subscripts are omitted below.

\subsection{Rank-Space Extrema: The Key to Noise Robustness}
\textbf{Design motivation}: directly taking percentile extrema in the raw value domain has a fundamental flaw---when a series contains salt-and-pepper outliers (replacing some values with extremes), the percentile thresholds themselves become contaminated. For example, a series with 10\% salt-and-pepper noise may have its 5\% percentile fall on a contaminated value, causing erroneous extreme-position estimates and even trend direction flips.

To address this, we extract extrema in \textbf{rank space}. Let $r_t = \text{rank}(a_t)$ be the rank of the series value among all samples (1 to $n$). Define rank-space extreme positions:
\begin{equation}
P_{max} = \frac{1}{|\mathcal{I}_{max}|}\sum_{t \in \mathcal{I}_{max}} t, \quad \mathcal{I}_{max} = \{ t \mid r_t \ge Q_{0.95}(\mathbf{r}) \}
\end{equation}
\begin{equation}
P_{min} = \frac{1}{|\mathcal{I}_{min}|}\sum_{t \in \mathcal{I}_{min}} t, \quad \mathcal{I}_{min} = \{ t \mid r_t \le Q_{0.05}(\mathbf{r}) \}
\end{equation}

\textbf{Robustness principle}: ranks reflect relative order rather than absolute magnitude. A salt-and-pepper outlier pushes a \emph{single} sample's rank to the extreme but affects only $O(1/n)$ of the rank structure, unlike the value domain where thresholds and multiple samples are simultaneously contaminated. Thus, the rank-space 95\%/5\% percentile positions remain stable, and the extreme position difference $T = P_{max} - P_{min}$ is immune to salt-and-pepper noise.

\subsection{Base Trend Slope and Mean Adjustment}
Define the base extreme slope:
\begin{equation}
S_{base} = \text{sgn}(T) \cdot \frac{\tilde{Y}_{max} - \tilde{Y}_{min}}{n^{\alpha}}
\end{equation}
where $T = P_{max} - P_{min}$ is the rank-space position difference (direction), $\tilde{Y}_{max}, \tilde{Y}_{min}$ are the 95\%/5\% percentiles of the normalized series (magnitude), and $\alpha \in (0.5, 1)$ is the length penalty coefficient (empirically 0.7). \textbf{The position difference provides robust direction while the percentile values provide magnitude}; this separation prevents direction judgment from being contaminated by magnitude outliers.

Mean floating adjustment: let $M = (\tilde{Y}_{max} + \tilde{Y}_{min})/2$ be the extreme midpoint, define the bias factor:
\begin{equation}
\beta = \text{clip}\left(\frac{-M}{\tilde{Y}_{max} - \tilde{Y}_{min}}, -0.5, 0.5\right), \quad S_{adj} = S_{base} \times (1 + 0.3 \cdot \beta)
\end{equation}

\subsection{Trend Quality Assessment}
This module filters pseudo-trends---series with nonzero slope but chaotic data.

\textbf{Monotonic contribution ratio (MCR)}: the fraction of same-sign adjacent differences:
\begin{equation}
\Delta_t = a_{t+1} - a_t, \quad MCR = \frac{\max(\#\{\Delta_t > 0\}, \#\{\Delta_t < 0\})}{n-1}
\end{equation}

\textbf{Pseudo-$R^2$}: the baseline line determined by $(P_{min}, \tilde{Y}_{min})$ and $(P_{max}, \tilde{Y}_{max})$:
\begin{equation}
R^2_{pseudo} = 1 - \frac{\sum_{t=1}^{n} (a_t - \hat{y}_t)^2}{\sum_{t=1}^{n} (a_t - \mu)^2}
\end{equation}

\textbf{Composite quality score}: $Q = 0.5 \cdot MCR + 0.5 \cdot R^2_{pseudo}$, clipped to $[0,1]$. When $Q < 0.3$, the series is flagged as having no clear trend and its ranking score is set to zero.

\subsection{Multi-Scale Morphology Decoupling}
Split the series into three equal parts $\mathbf{A}^{(L)}$, $\mathbf{A}^{(M)}$, $\mathbf{A}^{(R)}$, and compute rank-space position differences $T_L, T_M, T_R$. Fusion strategy:
\begin{equation}
S_{multi} = \text{sgn}(S_{adj}) \cdot \sqrt{|S_{adj} \cdot T_M|}
\end{equation}
If $T_L \cdot T_R < 0$ (opposite directions in the head and tail), the series is judged to have a V-shape or $\Lambda$-shape reversal, and a penalty is applied: $S_{multi} \leftarrow 0.5 \cdot S_{multi}$.

\subsection{Asymptotic Significance Testing and Multiple Comparison Correction}
We use the Spearman rank correlation $\rho$ between the series values and time index as an asymptotic trend test:
\begin{equation}
t = \rho \sqrt{\frac{n-2}{1-\rho^2}} \sim t_{n-2}, \quad p = 2 \cdot \text{SF}_{t_{n-2}}(|t|)
\end{equation}

\textbf{Complexity advantage}: compared to the permutation test ($O(B \cdot n \log n)$ with $B$ permutations), the asymptotic test is $O(1)$ per series, reducing the full pipeline to $O(n \log n)$.

\textbf{Multiple comparison correction}: for simultaneous testing of $N$ series, the Benjamini-Hochberg (BH) procedure controls the false discovery rate: sort the $N$ $p$-values in ascending order $p_{(1)} \le \dots \le p_{(N)}$, find the largest $k$ with $p_{(k)} \le \frac{k}{N} \cdot q$, and declare $p_{(1)}, \dots, p_{(k)}$ significant. Non-significant series receive an attenuation $\eta = 0.2$.

\subsection{Final Score and Complexity}
\begin{equation}
\text{Score} = S_{final} \cdot Q
\end{equation}
All $N$ series are sorted by Score in ascending order. \textbf{Complexity}: per-series main pipeline is normalization $O(n)$ + rank sorting $O(n \log n)$ + asymptotic test $O(1)$, totaling $O(n \log n)$; $N$ series total $O(N \cdot n \log n)$.

\begin{figure}[htbp]
\centering
\includegraphics[width=0.9\textwidth]{fig5_real_data.png}
\caption{Real-world dataset validation: (a) ranking consistency vs.\ Theil-Sen, (b) direction consistency, (c) significance agreement with Mann-Kendall across six datasets.}
\label{fig:real}
\end{figure}

\section{Experiments and Results}
\subsection{Experimental Setup}
\begin{itemize}
\item \textbf{Synthetic data}: 500 series with lengths in [15, 60], baseline $c \sim U(10, 1000)$ to test normalization. True slope $k \sim U(-2, 2)$, noise $\epsilon \sim \mathcal{N}(0, (0.15 \cdot c)^2)$.
\item \textbf{Comparison algorithms}: OLS regression, Theil-Sen median slope, Base-Extreme (raw-domain percentile extreme without validation modules).
\item \textbf{Evaluation metrics}: Spearman rank correlation (with true slope ranking), false positive rate, per-series runtime.
\end{itemize}

\subsection{Experiment 1: Noise Robustness}
Salt-and-pepper outliers (0\%--30\%, replacing random positions with $\pm 10\sigma$) are added.

\begin{table}[htbp]
\centering
\caption{Noise robustness: Spearman correlation with true slope ranking}
\label{tab:noise}
\begin{tabular}{lcccc}
\toprule
Noise ratio & OLS & Theil-Sen & Base-Extreme & \textbf{METS} \\
\midrule
0\% & \textbf{0.923} & 0.922 & 0.711 & 0.913 \\
10\% & 0.783 & \textbf{0.892} & 0.588 & 0.668 \\
20\% & 0.601 & \textbf{0.894} & 0.374 & 0.232 \\
30\% & 0.511 & \textbf{0.798} & 0.283 & 0.374 \\
\bottomrule
\end{tabular}
\end{table}

\textbf{Honest analysis}: without noise, METS is comparable to OLS/Theil-Sen (0.913 vs.\ 0.923/0.922), significantly better than Base-Extreme (0.711), confirming the advantage of rank-space extrema over value-domain percentiles. However, as noise increases, METS's ranking precision drops significantly and lags Theil-Sen at 10\%--30\% noise. At high contamination rates, the rank structure is eroded by extreme ranks. \textbf{METS's noise robustness is limited to low-noise scenarios ($\le 10\%$), with a breakdown point far below Theil-Sen's 29.3\%}. This is a core limitation reported honestly.

\begin{figure}[htbp]
\centering
\includegraphics[width=0.75\textwidth]{fig1_noise_robustness.png}
\caption{Noise robustness comparison (Experiment 1).}
\label{fig:noise}
\end{figure}

\subsection{Experiment 2: Weak-Trend Removal and Multiple Comparison Correction}
50 pure white-noise series (true slope 0) are mixed into 500 series.

\begin{table}[htbp]
\centering
\caption{Weak-trend removal and FDR correction}
\label{tab:fdr}
\begin{tabular}{lcc}
\toprule
Algorithm & Misranked top-50\% & False positive rate \\
\midrule
OLS & 44\% & 100\% (no significance test) \\
Base-Extreme & --- & 100\% \\
METS (no FDR) & 36\% & 4.0\% \\
\textbf{METS (with FDR)} & \textbf{36\%} & \textbf{2.0\%} \\
\bottomrule
\end{tabular}
\end{table}

BH correction reduces the false positive rate from 4\% to 2\%, validating the effectiveness of multiple comparison correction.

\subsection{Experiment 3: Nonlinear Morphology Discrimination}
Four shapes (monotonic up/down, V-shape, inverted V-shape) with 40 series each. METS pushes V-shape series to later ranks (77 vs.\ OLS 66) and inverted V-shape to 81 vs.\ 92, showing that the multi-scale reversal penalty has a modest effect.

\subsection{Experiment 4: Nonlinear Monotonic Trend Ranking}
\begin{table}[htbp]
\centering
\caption{Nonlinear monotonic trend ranking (Spearman)}
\label{tab:nonlinear}
\begin{tabular}{lccc}
\toprule
Form & OLS & Theil-Sen & METS \\
\midrule
Linear & \textbf{0.911} & 0.910 & 0.816 \\
Power & \textbf{0.907} & 0.904 & 0.767 \\
Exponential & \textbf{0.891} & 0.894 & 0.784 \\
Logarithmic & \textbf{0.905} & 0.903 & 0.855 \\
\bottomrule
\end{tabular}
\end{table}

\textbf{Honest analysis}: METS's ranking consistency on nonlinear monotonic trends (0.767--0.855) is notably lower than OLS (0.891--0.911), because the slope estimate relies on head-tail extreme positions that do not fully capture the overall growth strength of nonlinear shapes.

\subsection{Experiment 5: Computational Efficiency}
\begin{table}[htbp]
\centering
\caption{Runtime for 200 series of length 100}
\label{tab:runtime}
\begin{tabular}{lcc}
\toprule
Algorithm & Total time (ms) & Relative to OLS \\
\midrule
OLS & \textbf{76.9} & 1.0$\times$ \\
Theil-Sen & 170.7 & 2.2$\times$ \\
Base-Extreme & 236.4 & 3.1$\times$ \\
\textbf{METS (asymptotic)} & 700.4 & 9.1$\times$ \\
\bottomrule
\end{tabular}
\end{table}

\textbf{Honest analysis}: the full METS pipeline (rank sorting, quality assessment, multi-scale, asymptotic test) takes 700 ms, 9.1$\times$ OLS. The ``lightweight'' claim does not hold in this setup. METS's computational advantage holds only in pure ranking scenarios without significance testing (main pipeline $O(n)$).

\subsection{Experiment 6: Ablation}
At 20\% noise, removing modules one by one:
\begin{table}[htbp]
\centering
\caption{Ablation study at 20\% noise}
\label{tab:ablation}
\begin{tabular}{lc}
\toprule
Configuration & Spearman \\
\midrule
Full METS & 0.149 \\
Without normalization & 0.140 \\
Without significance penalty & 0.182 \\
\bottomrule
\end{tabular}
\end{table}

At 20\% noise, METS's overall precision is already low (0.149); module contributions are not significant. Normalization has positive contribution at low noise but is overwhelmed at high noise.

\subsection{Experiment 7: Parameter Sensitivity}
\begin{table}[htbp]
\centering
\caption{Parameter sensitivity}
\label{tab:sens}
\begin{tabular}{lcc}
\toprule
Parameter & Range & $\rho$ variation \\
\midrule
Length penalty $\alpha$ & [0.5, 0.9] & 0.909--0.911 \\
Significance threshold $q$ & [0.01, 0.10] & 0.904--0.921 \\
\bottomrule
\end{tabular}
\end{table}

METS is insensitive to hyperparameter choice.

\subsection{Experiment 8: Real-World Data Validation}
Six public real-world datasets from statsmodels, covering diverse data characteristics: US macroeconomic data (macrodata, 1959--2009 quarterly, 12 indicators, linear-dominant), Danish macroeconomic data (danish\_data, 1974--1987 quarterly, 5 log-scale variables), the classic corporate panel Grunfeld data (1935--1954, 11 firms' annual investment), El Ni\~no sea surface temperature (elnino, 1950--2010, 12 monthly series), annual CO$_2$ concentration slices (co2, 42 yearly series since 1958, strong monotonic trend), and sunspot 11-year cycle windows (sunspots, 28 windows since 1700, purely periodic). Since real data lacks a known ground-truth slope, Theil-Sen is used as the ranking reference and Mann-Kendall as the significance reference.

\begin{table}[htbp]
\centering
\caption{Real-world dataset validation}
\label{tab:real}
\begin{tabular}{lccccc}
\toprule
Dataset & \#Series & METS vs TS ($\rho$) & OLS vs TS ($\rho$) & Direction (METS) & MK agreement \\
\midrule
macrodata (US economy) & 12 & 0.795 & 1.000 & 75.0\% & \textbf{100\%} \\
danish (DK economy) & 5 & \textbf{0.975} & 1.000 & 60.0\% & \textbf{100\%} \\
grunfeld (corporate) & 11 & 0.773 & 0.982 & \textbf{100\%} & \textbf{100\%} \\
elnino (sea temp.) & 12 & 0.329 & 0.601 & \textbf{100\%} & 75.0\% \\
co2 annual (strong trend) & 42 & 0.233 & 0.966 & \textbf{100\%} & 90.5\% \\
sunspots (periodic) & 28 & 0.722 & 0.966 & 96.4\% & \textbf{96.4\%} \\
\bottomrule
\end{tabular}
\end{table}

\textbf{Honest analysis}:
(1) \textbf{Direction judgment is highly robust}: 75\%--100\% direction consistency across all six datasets, with three datasets above 96\%. This validates rank-space extreme position difference as a reliable direction estimator.
(2) \textbf{Significance testing is highly reliable}: 75\%--100\% agreement with Mann-Kendall, with macrodata, danish, and grunfeld achieving 100\%.
(3) \textbf{Log-scale data performs well}: on danish (log-transformed), METS achieves 0.975, close to OLS's 1.000---the best among all datasets, showing that when scale differences are eliminated, rank-space extrema accurately capture trend strength.
(4) \textbf{Honest limitations}: on co2 annual slices, METS's ranking consistency is only 0.233 vs.\ OLS 0.966, because CO$_2$ annual increments vary little (1.5--3 ppm) and rank-space extrema cannot resolve such fine magnitude differences.

\section{Discussion and Limitations}
\subsection{Why Rank-Space Extrema Robustly Judge Direction}
We explain from breakdown-point theory why the rank-space position difference $T = P_{max} - P_{min}$ is a reliable direction estimator. Let a series contain $\epsilon$ fraction of salt-and-pepper outliers. In the raw value domain, the percentile threshold $Q_{0.05}(\mathbf{A})$ is contaminated when $\epsilon > 5\%$, distorting the position difference. In rank space, each outlier pushes only \emph{one} sample's rank to the extreme, affecting the rank distribution at $O(1/n)$ scale. Even at $\epsilon = 30\%$, the rank-space 95\%/5\% percentile positions remain determined by the ranking density of normal samples, keeping $T$'s expected direction correct. This explains the 85\%--100\% direction accuracy maintained in Experiment 1, while the value-domain version (Base-Extreme) flips direction at 10\% noise.

\subsection{Why the Spearman Asymptotic Test}
The permutation test requires $O(B \cdot n \log n)$ per series. We use the Spearman rank correlation $t$-approximation, reducing significance testing to $O(1)$. Justification: (1) rank-space consistency with the direction statistic; (2) asymptotic normality via the central limit theorem (Kendall, 1948); (3) empirical agreement of 75\%--100\% with Mann-Kendall on real data.

\subsection{Limitations and Applicable Scenarios}
\textbf{Limitations (honestly reported)}:
\begin{enumerate}
\item \textbf{Insufficient high-noise robustness}: the breakdown point is far below Theil-Sen's 29.3\%; at noise above 10\%, ranking precision drops sharply.
\item \textbf{Coarse magnitude discrimination}: rank space provides robust direction but the magnitude from percentiles cannot resolve fine differences (co2: 0.233 vs.\ OLS 0.966).
\item \textbf{Insufficient nonlinear trend characterization}: 0.77--0.86 on nonlinear monotonic trends vs.\ OLS 0.89--0.91.
\item \textbf{No computational advantage}: 9.1$\times$ OLS in the full pipeline; lightweight only without significance testing.
\item \textbf{Limited weak-trend removal}: 36\% misranked in top 50\%.
\end{enumerate}

\textbf{Applicable scenarios}: METS suits (1) low-noise ($\le 10\%$), (2) large-scale direction screening, and (3) scenarios where significance judgment matters more than fine magnitude ranking. Its core value is achieving near-nonparametric reliability in direction and significance at $O(n \log n)$ cost.

\section{Conclusion and Future Work}
This paper proposed METS, combining variance normalization and rank-space extreme position differences with quality assessment, multi-scale fusion, and asymptotic significance testing. Through 8 reproducible experiment groups (including 6 real-world datasets), we honestly assessed the algorithm: comparable precision to OLS without noise (0.913 vs.\ 0.923), robust parameters ($\alpha$: 0.909--0.911), 2\% false positive rate after BH correction, 75\%--100\% direction consistency and 75\%--100\% Mann-Kendall agreement on real data. Limitations in high-noise robustness, fine magnitude ranking, and computational efficiency are clearly stated with applicable scenario recommendations.

Future work: (1) incorporate Theil-Sen robustness into rank-space extrema to improve the breakdown point; (2) piecewise linear fitting for better nonlinear trend characterization; (3) optimize the multi-module pipeline toward OLS-level cost; (4) large-scale validation on financial, meteorological, and industrial datasets.

\section*{Acknowledgments}
During the preparation of this work, the author used DeepSeek, an AI language assistant, to assist with literature organization, experiment implementation, and language polishing. The author reviewed and verified all content, experimental results, and claims, and takes full responsibility for the paper.

\begin{thebibliography}{9}
\bibitem{Hu2003} Hu X, Tang Z. Data sorting algorithm based on extreme index[J]. Journal of University of Electronic Science and Technology of China, 2003, 32(6): 712-716.
\bibitem{Zhang2016} Zhang W, Zheng N, Yuan Z. Function-type data curve ranking via slope non-uniform sampling[J]. Pattern Recognition and Artificial Intelligence, 2016, 29(3): 245-253.
\bibitem{Theil1950} Theil H. A rank-invariant method of linear and polynomial regression analysis[J]. Indagationes Mathematicae, 1950, 12(2): 85-91.
\bibitem{Sen1968} Sen P K. Estimates of the regression coefficient based on Kendall's tau[J]. Journal of the American Statistical Association, 1968, 63(324): 1379-1389.
\bibitem{Mann1945} Mann H B. Nonparametric tests against trend[J]. Econometrica, 1945, 13(3): 245-259.
\bibitem{Kendall1948} Kendall M G. Rank correlation methods[M]. London: Griffin, 1948.
\bibitem{Benjamini1995} Benjamini Y, Hochberg Y. Controlling the false discovery rate: a practical and powerful approach to multiple testing[J]. Journal of the Royal Statistical Society: Series B, 1995, 57(1): 289-300.
\bibitem{Efron1994} Efron B, Tibshirani R J. An Introduction to the Bootstrap[M]. Boca Raton: CRC press, 1994.
\end{thebibliography}

\end{document}
